\documentclass[11pt,a4paper]{article}

\newcommand{\rapporttitre}{Prévision ponctuelle}
\newcommand{\rapportversion}{v0.1}
% Mise en page commune à tous les rapports du projet.
% Chaque rapport définit \rapporttitre et \rapportversion avant de charger ce fichier.

\usepackage{fontspec}
\usepackage[french]{babel}
\usepackage[a4paper,margin=2.2cm,bottom=2.6cm]{geometry}
\usepackage{graphicx}
\usepackage{float}
\usepackage{booktabs}
\usepackage{tabularx}
\usepackage{array}
\usepackage[table]{xcolor}
\usepackage{caption}
\usepackage{enumitem}
\usepackage{titlesec}
\usepackage{fancyhdr}
\usepackage{hyperref}
\usepackage{xurl}

\setmainfont{texgyreheros}[
  Extension=.otf,
  UprightFont=*-regular,
  BoldFont=*-bold,
  ItalicFont=*-italic,
  BoldItalicFont=*-bolditalic]
\definecolor{accent}{HTML}{2A78D6}
\definecolor{ink2}{HTML}{52514E}
\definecolor{fond}{HTML}{EEF4FC}
\hypersetup{colorlinks=true, linkcolor=accent, urlcolor=accent}
\graphicspath{{figures/}}

\titleformat{\section}{\Large\bfseries\color{accent}}{\thesection.}{0.6em}{}[\vspace{-0.4em}\color{accent}\rule{\linewidth}{0.8pt}]
\titleformat{\subsection}{\large\bfseries}{\thesubsection}{0.6em}{}
\titlespacing*{\section}{0pt}{1.6em}{0.9em}

\captionsetup{font={small,color=ink2}, labelfont=bf, justification=raggedright, singlelinecheck=false}
\setlist{itemsep=0.25em, topsep=0.4em}
\setlength{\parskip}{0.45em}
\renewcommand{\arraystretch}{1.25}
\newcolumntype{Y}{>{\raggedright\arraybackslash}X}

\pagestyle{fancy}
\fancyhf{}
\fancyhead[L]{\small\color{ink2} day-ahead-load -- \rapporttitre}
\fancyhead[R]{\small\color{ink2} \rapportversion}
\fancyfoot[R]{\small\color{ink2}\thepage}
\renewcommand{\headrulewidth}{0.3pt}

% \pagedegarde{sous-titre}{date}{ligne de version}
\newcommand{\pagedegarde}[3]{%
  \hypersetup{pageanchor=false}
  \begin{titlepage}
    \setlength{\parindent}{0pt}
    \vspace*{5cm}
    {\color{accent}\bfseries\large DAY-AHEAD-LOAD}\par\vspace{0.6cm}
    {\fontsize{30}{36}\selectfont\bfseries \rapporttitre}\par\vspace{0.5cm}
    {\Large\color{ink2} #1}\par
    \vfill
    {\color{accent}\rule{\linewidth}{1.2pt}}\par\vspace{0.3cm}
    {\color{ink2} Nathan Mang\\ #2\\ #3}
    \vspace{2cm}
  \end{titlepage}
  \hypersetup{pageanchor=true}}

% Encadré pour les points à retenir.
\newcommand{\aretenir}[1]{%
  \par\noindent\colorbox{fond}{\parbox{\dimexpr\linewidth-2\fboxsep}{\small #1}}\par}

\usepackage{needspace}
\usepackage{etoolbox}
% Laisse un peu plus d'espace entre les mots plutôt que de déborder.
\setlength{\emergencystretch}{3em}
% Un titre de section ne reste jamais seul en bas de page.
\pretocmd{\section}{\needspace{10\baselineskip}}{}{}

\begin{document}

\pagedegarde{Choisir et mesurer la prévision de la consommation du lendemain, sur la validation seulement}{24 novembre 2026}{Version 0.1, test fermé}

\setcounter{tocdepth}{1}
\tableofcontents
\newpage

\section*{Synthèse}
\addcontentsline{toc}{section}{Synthèse}

La v0.1 choisit le modèle qui prévoit la consommation française du lendemain, par demi-heure, à partir des seules données publiées avant midi la veille. Le choix se fait sur la validation, du 15/03/2024 au 31/08/2025. Le test, du 01/09/2025 au 31/08/2026, reste fermé jusqu'à la v1.0.

\begin{itemize}
  \item Modèle retenu : LightGBM aux réglages par défaut, avec le niveau des 52 dernières semaines en variable, entraîné de 2017 à mars 2024 sur la météo ERA5. MAPE de validation : 2,61~\%.
  \item Chaque choix a suivi une règle écrite avant le calcul. Dix modèles ont été comparés sur la validation ; ce nombre est publié, parce qu'il rend le score de validation optimiste.
  \item Un modèle plus complexe n'est retenu que si son gain dépasse 1~\% et sort du bruit des semaines. Ce test a écarté une correction de la météo qui gagnait 1,6~\%.
  \item Le modèle bat tous les naïfs, ne se distingue pas de la prévision brute de RTE, publiée après la décision, et reste derrière ses versions corrigées.
  \item Avec une météo parfaite, l'erreur ne baisse que de 0,03 point : à J-1, l'erreur vient surtout du modèle.
\end{itemize}

\aretenir{Rapport au naïf J-7 sur la validation : 0,43, intervalle à 95~\% de 0,38 à 0,50. Le seuil de réussite fixé en v0.0 est 0,5. C'est une indication, pas le critère : ce chiffre est probablement optimiste, et le critère se juge sur le test.}

\section{Le cadre}

Un fournisseur décide à midi, heure de Paris, combien acheter pour chaque demi-heure du lendemain. Le cadre a été fixé en v0.0, avant toute modélisation (ADR 0001) : découpage dans le temps, références, critères de réussite. La table de données est fixée par l'ADR 0002 :

\begin{itemize}
  \item une ligne par demi-heure réelle, indexée en UTC, avec 46 ou 50 demi-heures les jours de changement d'heure ;
  \item une cible consolidée ou définitive, publiée par RTE sur éCO2mix ;
  \item un calendrier : fériés, lundi de Pentecôte, ponts, période de Noël ;
  \item une météo nationale (température, nébulosité, vent, rayonnement) pondérée par la consommation régionale aux 12 points de la v0.0 ;
  \item le niveau de la consommation, moyenne des 52 dernières semaines de consolidé publié à midi la veille.
\end{itemize}

La règle la plus importante porte sur la météo. À l'entraînement, elle vient d'ERA5. En validation, en test et en production, elle vient du calcul ECMWF de 00~UTC de la veille, y compris dans les températures lissées sur 24~h et sur 72~h. Un test le vérifie : il remplace par une valeur absurde tout calcul postérieur à midi la veille, et les variables ne doivent pas changer.

\section{Une méthode qui protège la validation}

Chaque fois qu'on choisit un modèle en regardant la validation, son score sur cette période devient un peu plus flatteur. Trois règles limitent cet effet et le rendent visible.

\begin{itemize}
  \item \textbf{La règle avant le calcul.} Chaque décision est écrite et commitée dans le journal des choix avant le calcul qui la tranche. L'historique Git en garde la preuve.
  \item \textbf{Le compteur.} Chaque modèle évalué sur la validation a une entrée datée, avec le total cumulé. Il est de 10 à la fin de la v0.1. Le journal ne fait que s'allonger : une décision qu'on revoit est remplacée par une nouvelle entrée, jamais réécrite.
  \item \textbf{Le test de bruit.} Sur 17,5 mois, une différence de 1~\% peut venir du hasard de quelques semaines. Pour chaque comparaison, on calcule la différence appariée de MAPE et son intervalle à 95~\%, par un bootstrap qui tire des blocs d'une semaine (77 blocs, 10\,000 tirages). Le modèle plus complexe n'est retenu que s'il gagne plus de 1~\% \textbf{et} si l'intervalle exclut 0. Sinon, on garde le plus simple.
\end{itemize}

Le modèle de comparaison est figé : LightGBM 4.7.0, réglages par défaut, graine fixée, exécution déterministe. Seules les entrées changent d'une comparaison à l'autre. Chaque calcul relancé redonne un fichier identique.

\section{Les choix}

\begin{table}[H]
\small
\begin{tabularx}{\linewidth}{lYl}
\toprule
\textbf{Décision} & \textbf{Mesure} & \textbf{Retenu} \\
\midrule
Convention de la demi-heure & $-2{,}1$~\% en validation croisée 2019-2023, dans l'entraînement & $(t-15, t+15]$ \\
Source météo & ERA5 corrigé : $+1{,}6$~\%, intervalle $[-0{,}20 ; +0{,}12]$ ; IFS : $+0{,}5$~\% & ERA5 \\
Durée & 2017 contre 2012 : $+5{,}8$~\%, intervalle $[-0{,}27 ; -0{,}05]$ & depuis 2017 \\
Famille et niveau & LightGBM contre linéaire, avec niveau : $+28{,}3$~\%, intervalle $[-1{,}38 ; -0{,}68]$ & LightGBM, niveau \\
\bottomrule
\end{tabularx}
\caption{Les intervalles portent sur la différence de MAPE, en points. Un gain dont l'intervalle contient 0 est traité comme du bruit.}
\end{table}

\subsection{Convention de la demi-heure}

La consommation étiquetée $t$ est une puissance moyenne, mais la documentation ne dit pas sur quel intervalle. Deux lectures ont été comparées à l'intérieur de l'entraînement, sans toucher à la validation, en mettant de côté chaque année de 2019 à 2023. L'intervalle centré $(t-15, t+15]$ gagne 2,1~\%, au-delà du seuil de 1~\% fixé d'avance. Le signal est fragile : il gagne trois années sur cinq.

\subsection{Source de la météo d'entraînement}

ERA5 est une réanalyse ; la météo prévue vient d'un autre modèle. Trois sources d'entraînement ont été comparées : ERA5 brut, la série IFS historique (proche de la prévision, mais seulement depuis 2017) et ERA5 corrigé de l'écart moyen IFS~$-$~ERA5 par point, mois et heure. La version corrigée gagne 1,6~\%. Le seuil de 1~\% seul l'aurait retenue, mais l'intervalle de la différence contient 0 : le gain est de l'ordre du hasard des semaines, et ERA5 brut est gardé.

\subsection{Durée de l'entraînement}

Entraîné à partir de 2017, le modèle fait 5,8~\% mieux qu'à partir de 2012, et l'intervalle exclut 0. Les années 2012-2016 dégradent la prévision de 2024-2025. Deux explications restent possibles, sans que la mesure les sépare : un changement des usages de l'électricité, ou l'absence du niveau avant mars 2013.

\subsection{Famille de modèles et niveau}

Six modèles ont été comparés : une régression linéaire et LightGBM, chacun sans niveau, avec le niveau en variable, ou sur le rapport consommation~/~niveau. La régression linéaire n'est pas une référence de façade. Elle croise la température avec l'heure et avec le mois, comme la référence linéaire de la compétition GEFCom2012 : l'effet du chauffage n'est pas le même à 3~h et à 19~h.

\begin{table}[H]
\small
\begin{tabularx}{\linewidth}{Yrrr}
\toprule
\textbf{Famille} & \textbf{Sans niveau} & \textbf{Niveau en variable} & \textbf{Rapport} \\
\midrule
Régression linéaire & 4,18~\% & 3,65~\% & 3,85~\% \\
LightGBM & 4,35~\% & \textbf{2,62~\%} & 3,52~\% \\
\bottomrule
\end{tabularx}
\caption{MAPE de validation des six modèles. Les modèles sont parcourus du plus simple au plus complexe, chacun opposé au modèle retenu jusque-là.}
\end{table}

Le niveau pèse lourd : sans lui, LightGBM ne sait pas suivre la baisse de la consommation et fait moins bien que la régression linéaire. Le rapport consommation~/~niveau, censé corriger ce défaut, est moins bon dans les deux familles.

\section{Ce que vaut le modèle}

Le bilan mesure le modèle retenu sans rien choisir : aucune de ses mesures ne le modifie. Il porte sur les 25\,607 demi-heures de validation où le modèle et les dix références de la v0.0 existent toutes.

\begin{table}[H]
\small
\begin{tabularx}{\linewidth}{Ylrr}
\toprule
\textbf{Référence} & \textbf{Étiquette} & \textbf{MAPE} & \textbf{Modèle $-$ référence} \\
\midrule
Modèle retenu & & 2,61~\% & \\
Même demi-heure 7 jours avant (J-7) & optimiste & 6,04~\% & $-3{,}42$ $[-4{,}30 ; -2{,}59]$ \\
Même jour, 52 semaines avant & réaliste & 6,34~\% & $-3{,}73$ $[-4{,}60 ; -2{,}89]$ \\
Dernier même jour publié en consolidé & réaliste & 18,76~\% & $-16{,}14$ $[-19{,}01 ; -13{,}36]$ \\
RTE J-1 brute & avantagée & 2,58~\% & $+0{,}03$ $[-0{,}24 ; +0{,}31]$ \\
RTE J-1 corrigée (28 derniers jours publiés) & avantagée & 1,92~\% & $+0{,}70$ $[+0{,}45 ; +0{,}96]$ \\
RTE J-1 corrigée (même mois l'année précédente) & avantagée & 1,72~\% & $+0{,}89$ $[+0{,}68 ; +1{,}13]$ \\
\bottomrule
\end{tabularx}
\caption{Différences en points de MAPE, avec leur intervalle à 95~\%. « Avantagée » : publiée après la décision. Les autres références de la v0.0 sont dans le fichier de résultats.}
\end{table}

Face à J-7, le rapport des MAPE est de 0,43, avec un intervalle de 0,38 à 0,50. Ce chiffre est probablement optimiste, pour trois raisons : dix comparaisons ont été faites sur cette période, la météo prévue de la validation est rejouée (cycle 49R1 d'ECMWF), et J-7 lit du consolidé qu'un fournisseur n'aurait pas encore à midi. Le critère de l'ADR 0001 se juge sur le test.

\subsection{La part de la météo prévue}

Pour savoir ce que coûte l'erreur de la prévision météo, le même modèle a prévu la validation avec la météo observée (ERA5) en entrée. Cette borne est impossible en pratique : la météo du lendemain n'est jamais connue. La MAPE passe de 2,614~\% à 2,583~\%, soit $-0{,}03$ point, avec un intervalle de $-0{,}17$ à $+0{,}09$ qui contient 0.

\aretenir{À J-1, la prévision météo est assez bonne pour que son erreur pèse peu. L'erreur à couvrir par l'incertitude de la v0.2 vient surtout du modèle.}

\section{Où le modèle se trompe}

Le biais est mesuré comme la prévision moins la consommation : positif, le modèle prévoit trop. En moyenne, il prévoit trop d'environ 210~MW.

\begin{table}[H]
\small
\begin{tabularx}{\linewidth}{Yrrr}
\toprule
\textbf{Température de l'indice (ERA5)} & \textbf{Demi-heures} & \textbf{MAPE} & \textbf{Biais} \\
\midrule
Sous 0~°C & 214 & 2,03~\% & $+845$~MW \\
De 0 à 5~°C & 1\,975 & 2,47~\% & $+764$~MW \\
De 5 à 10~°C & 4\,727 & 3,19~\% & $+673$~MW \\
De 10 à 15~°C & 6\,290 & 2,89~\% & $+455$~MW \\
Au-dessus de 15~°C & 12\,401 & 2,29~\% & $-177$~MW \\
\bottomrule
\end{tabularx}
\end{table}

Par temps froid, le modèle prévoit trop de 0,67 à 0,85~GW. C'est le signe et l'ordre de grandeur du risque de 720~MW noté au prémortem de la v0.0, qui venait de l'écart entre ERA5 et la météo prévue. Mais avec la météo parfaite, l'erreur ne change presque pas : la source météo n'explique donc pas ce biais, dont la cause n'est pas établie.

Autres constats :

\begin{itemize}
  \item La pire tranche horaire est 6~h--9~h (MAPE 3,34~\%, $+940$~MW), puis 3~h--6~h. Entre 12~h et 15~h, le modèle prévoit trop peu ($-455$~MW).
  \item Les dimanches et fériés sont moins bien prévus que les jours ouvrés (3,14~\% contre 2,45~\%). Les ponts ($+897$~MW) et la période de Noël ($+1\,508$~MW) sont nettement surestimés : ces jours sont rares à l'entraînement.
  \item Par trimestre, la MAPE va de 2,25~\% à 2,87~\%, sauf au premier trimestre 2024 (4,46~\%), expliqué plus bas.
\end{itemize}

\subsection{Les pires jours}

Pour les dix pires jours, une explication n'est écrite que si elle se vérifie dans nos données : calendrier, froid ou chaleur marqués, erreur de la météo prévue. Sinon, la cause n'est pas établie.

\begin{table}[H]
\small
\begin{tabularx}{\linewidth}{lrrY}
\toprule
\textbf{Jour} & \textbf{MAPE} & \textbf{Biais} & \textbf{Constat} \\
\midrule
16/03/2024 & 17,7~\% & $+8{,}2$~GW & non établi (voir ci-dessous) \\
15/03/2024 & 14,2~\% & $+7{,}2$~GW & non établi (voir ci-dessous) \\
02/01/2025 & 9,6~\% & $+5{,}9$~GW & non établi \\
09/04/2024 & 9,1~\% & $-4{,}5$~GW & non établi \\
15/09/2024 & 8,9~\% & $+3{,}4$~GW & froid marqué \\
25/12/2024 & 8,6~\% & $+4{,}4$~GW & férié, période de Noël \\
14/09/2024 & 7,6~\% & $+3{,}0$~GW & froid marqué \\
06/10/2024 & 7,5~\% & $+3{,}0$~GW & non établi \\
08/04/2024 & 7,3~\% & $-3{,}3$~GW & non établi \\
03/01/2025 & 7,2~\% & $+4{,}7$~GW & non établi \\
\bottomrule
\end{tabularx}
\end{table}

Les deux pires jours, les 15 et 16/03/2024, sont les premiers de la validation. Ce sont aussi exactement les 96 demi-heures où la température lissée sur 72~h manque : les calculs ECMWF en cache ne remontent pas assez loin pour couvrir leurs 72 heures. LightGBM envoie ces valeurs manquantes dans une branche qui surestime de 7 à 8~GW. En production, le calcul de secours aurait pris le relais. Ce constat a été fait après le calcul, hors de la règle des constats : il est noté comme une observation, et il gonfle la MAPE de validation d'environ 0,06 point.

\section{Les variables qui comptent}

\begin{table}[H]
\small
\begin{tabularx}{\linewidth}{Yrr}
\toprule
\textbf{Variable} & \textbf{Permutation (points de MAPE)} & \textbf{Gain (part)} \\
\midrule
Demi-heure & $+8{,}65$ & 14~\% \\
Température lissée sur 72~h & $+6{,}45$ & 53~\% \\
Jour de la semaine & $+3{,}90$ & 6~\% \\
Température lissée sur 24~h & $+2{,}38$ & 19~\% \\
Mois & $+1{,}36$ & 3~\% \\
Température de la demi-heure & $+0{,}73$ & 1~\% \\
Niveau des 52 semaines & $+0{,}14$ & 2~\% \\
\bottomrule
\end{tabularx}
\caption{Permutation : hausse de la MAPE de validation quand la variable est mélangée, moyenne de 5 tirages. Gain : part du gain total des coupures, sur l'entraînement.}
\end{table}

Les deux mesures ont leurs limites. Le gain favorise les variables qui offrent beaucoup de coupures possibles. La permutation sous-estime une variable liée à d'autres : mélanger la température laisse l'information dans ses lissages. Elle ne voit pas non plus le rôle du niveau, qui décale toute la validation par rapport à l'entraînement mais varie peu à l'intérieur. Avec le niveau, le modèle gagne pourtant 40~\% sur le même modèle sans niveau.

\section{Limites et points ouverts}

\begin{itemize}
  \item Le score de validation est optimiste : dix comparaisons, une météo rejouée, un critère jugé face à un naïf lui-même optimiste. Le test le mesurera.
  \item Faut-il interdire les valeurs d'entrée manquantes dans LightGBM, par le calcul de secours ou par un refus de prévoir ?
  \item D'où vient le biais par temps froid, alors que la météo parfaite ne change presque rien ?
  \item Pourquoi les années 2012-2016 dégradent-elles la prévision ?
  \item Les vacances scolaires et le régime local d'Alsace-Moselle ne sont pas dans le modèle.
  \item Le seuil de climatisation de la régression linéaire (18~°C) a été fixé sans mesure ; ERA5 corrigé ramène 19~\% des heures de nébulosité dans leurs bornes.
\end{itemize}

\section{Suite}

La v0.2 ajoute l'incertitude autour de ce modèle, sans le modifier : intervalles à 80~\% et à 95~\% et score quantile, comparés à la référence de quantiles de l'ADR 0001. Le bilan oriente ce travail : l'erreur à couvrir vient surtout du modèle, elle est plus forte le matin, les dimanches, les ponts et à Noël, et biaisée par temps froid.

Le détail de chaque choix est dans le journal des choix et dans l'ADR 0003. Les chiffres viennent des scripts de \nolinkurl{scripts/modelling}, dont les sorties sont dans \nolinkurl{docs/v0.1-prevision-ponctuelle/data}.

\end{document}
